\documentclass[12pt]{article}
\usepackage[utf8]{inputenc}
\usepackage[spanish]{babel}
\usepackage{amsmath}
\usepackage{amssymb}
\usepackage{geometry}
\usepackage{hyperref}
\usepackage{booktabs}
\usepackage{graphicx}

\geometry{margin=2.5cm}

\title{Tema 4: Fundamentación para la resolución del problema del ciclo hamiltoniano}
\author{}
\date{}

\begin{document}

\maketitle

\section{Introducción}

El problema del ciclo hamiltoniano (\textit{Hamiltonian Cycle Problem}, HCP) consiste en determinar si un grafo \(G=(V,E)\) contiene un ciclo que visite cada vértice exactamente una vez y regrese al vértice inicial. Este problema es fundamental en teoría de grafos, optimización combinatoria e inteligencia artificial, ya que pertenece a la clase de problemas NP-completos. \cite{fhcp}

El presente documento fundamenta el diseño de una solución basada en algoritmos metaheurísticos: Recocido Simulado (SA), Optimización por Colonia de Hormigas (ACO), Colonia de Abejas (ABC), Evolución Diferencial (DE) y Optimización por Enjambre de Partículas (PSO). La propuesta se evaluará utilizando el FHCP Challenge Set, un conjunto de 1{,}001 instancias estructuralmente difíciles del HCP. \cite{fhcpset}

\section{Teoría de grafos y ciclos hamiltonianos}

Un grafo se define formalmente como \(G=(V,E)\), donde \(V\) es el conjunto de vértices y \(E\) es el conjunto de aristas. En un grafo no dirigido, una arista \(\{u,v\}\) conecta dos vértices sin dirección; en un grafo dirigido, la arista \((u,v)\) tiene orientación de \(u\) hacia \(v\). \cite{west}

\subsection{Definiciones formales}

\begin{itemize}
    \item \textbf{Camino:} secuencia de vértices \(v_1,v_2,\dots,v_k\) tal que existe una arista entre cada par consecutivo y ningún vértice se repite.
    \item \textbf{Ciclo:} camino cerrado, es decir, una secuencia \(v_1,v_2,\dots,v_k,v_1\) donde los vértices \(v_1,\dots,v_k\) son distintos.
    \item \textbf{Ciclo hamiltoniano:} ciclo que visita todos los vértices del grafo exactamente una vez.
    \item \textbf{Grafo hamiltoniano:} grafo que contiene al menos un ciclo hamiltoniano.
\end{itemize}

Dado un grafo con \(n=|V|\) vértices, un ciclo hamiltoniano puede representarse como una permutación cíclica:

\[
\pi=(v_1,v_2,\dots,v_n)
\]

donde cada vértice aparece exactamente una vez y se cumple que:

\[
\{v_i,v_{i+1}\}\in E
\]

para todo \(i=1,\dots,n\), considerando que \(v_{n+1}=v_1\).

Por lo tanto, un ciclo hamiltoniano tiene longitud exactamente \(n\): utiliza \(n\) aristas para recorrer los \(n\) vértices y regresar al origen. \cite{fhcp}

\subsection{Diferencia entre camino y ciclo hamiltoniano}

\begin{table}[h]
\centering
\begin{tabular}{lll}
\toprule
\textbf{Concepto} & \textbf{Definición} & \textbf{Condición final} \\
\midrule
Camino hamiltoniano & Visita todos los vértices una sola vez & No requiere regresar al inicio \\
Ciclo hamiltoniano & Visita todos los vértices una sola vez & Debe regresar al vértice inicial \\
\bottomrule
\end{tabular}
\end{table}

La diferencia esencial es que el ciclo hamiltoniano exige una arista adicional que cierre la secuencia. El FHCP define el problema como encontrar un ciclo simple de longitud \(N\), donde \(N\) es el número total de vértices. \cite{fhcp}

\section{Representación del espacio de búsqueda}

La representación del problema determina qué soluciones puede explorar el algoritmo y qué operadores pueden aplicarse. En el HCP, una solución válida debe ser una permutación cíclica de todos los vértices; no puede contener repeticiones, omisiones ni aristas inexistentes.

\subsection{Codificación por permutación}

La representación más adecuada es una permutación de vértices:

\[
x=(v_{\pi(1)},v_{\pi(2)},\dots,v_{\pi(n)})
\]

donde \(\pi\) es una biyección sobre el conjunto de vértices. Esta codificación garantiza que cada vértice aparezca exactamente una vez.

Por ejemplo, para un grafo con vértices \(\{A,B,C,D\}\), la secuencia:

\[
(A,B,C,D,A)
\]

representa un ciclo candidato. Este ciclo es válido únicamente si las aristas \(AB\), \(BC\), \(CD\) y \(DA\) existen en el grafo.

\subsection{Codificaciones posibles}

\begin{table}[h]
\centering
\begin{tabular}{lll}
\toprule
\textbf{Codificación} & \textbf{Descripción} & \textbf{Riesgo principal} \\
\midrule
Permutación directa & Lista ordenada de los \(n\) vértices & Puede generar aristas inexistentes \\
Matriz de adyacencia binaria & Cada vértice selecciona dos vecinos & Puede producir subciclos disjuntos \\
Selección de aristas & Subconjunto de \(n\) aristas con grado 2 & Puede formar varios ciclos \\
Permutación con reparación & Se corrigen violaciones después de generar la solución & La reparación puede modificar la solución original \\
\bottomrule
\end{tabular}
\end{table}

Una selección de aristas donde cada vértice tiene grado 2 garantiza que todos los vértices participan en algún ciclo, pero no garantiza que todos pertenezcan al mismo ciclo. Por ello, además del grado 2, debe verificarse la conectividad del subgrafo seleccionado. \cite{fhcp}

\subsection{Operadores seguros para algoritmos evolutivos}

Para evitar generar soluciones inválidas, se recomienda utilizar operadores que preserven la estructura de permutación:

\begin{itemize}
    \item \textbf{Order Crossover (OX):} conserva el orden relativo de los vértices sin duplicarlos.
    \item \textbf{Partially Matched Crossover (PMX):} intercambia segmentos y repara duplicados.
    \item \textbf{Swap mutation:} intercambia dos posiciones de la permutación.
    \item \textbf{Inversion mutation:} invierte un segmento contiguo de la secuencia.
    \item \textbf{2-opt:} elimina dos aristas y las reconecta de otra forma.
\end{itemize}

La codificación debe evitar:

\begin{itemize}
    \item Vértices repetidos.
    \item Vértices omitidos.
    \item Aristas inexistentes.
    \item Ciclos parciales.
    \item Subciclos desconectados.
    \item Soluciones que no cierren correctamente.
\end{itemize}

\section{Módulo evolutivo y heurísticas}

Un algoritmo evolutivo representa cada solución candidata como un individuo, evalúa su calidad mediante una función de aptitud (\textit{fitness}) y genera nuevas soluciones mediante selección, cruce, mutación y reemplazo.

\subsection{Función de aptitud}

Una función de aptitud puede definirse como:

\[
f(x)=\alpha \cdot \text{violaciones}(x)+\beta \cdot \text{longitud}(x)
\]

donde:

\begin{itemize}
    \item \(\text{violaciones}(x)\) cuenta aristas inválidas, vértices repetidos, vértices omitidos o subciclos.
    \item \(\text{longitud}(x)\) mide el número de aristas utilizadas.
    \item \(\alpha\) y \(\beta\) son pesos que equilibran factibilidad y calidad.
\end{itemize}

Si el objetivo es únicamente encontrar un ciclo hamiltoniano válido, la función puede definirse de forma binaria:

\[
f(x)=
\begin{cases}
1, & \text{si } x \text{ es un ciclo hamiltoniano válido} \\
0, & \text{en otro caso}
\end{cases}
\]

Sin embargo, una función binaria no guía adecuadamente la búsqueda. Por ello, se recomienda una función gradual que premie soluciones parcialmente correctas.

\subsection{Heurísticas propuestas}

Una heurística es una regla que estima qué acción acerca más al algoritmo a la solución objetivo sin garantizar optimalidad. En inteligencia artificial, las heurísticas permiten reducir el espacio de búsqueda cuando la exploración exhaustiva es inviable. \cite{russell}

Para el HCP pueden utilizarse las siguientes heurísticas:

\begin{itemize}
    \item Selección de vértices de bajo grado.
    \item Evitar movimientos que fragmenten el grafo.
    \item Reparación de subciclos disjuntos.
    \item Aplicación de búsqueda local 2-opt.
    \item Construcción inicial mediante algoritmos greedy.
    \item Priorización de conexiones que mantengan la conectividad.
\end{itemize}

\section{Métricas de evaluación en inteligencia artificial}

La evaluación debe combinar métricas algorítmicas y métricas estadísticas. En un problema de decisión como el HCP, la métrica principal es determinar si el algoritmo encuentra un ciclo válido dentro de un presupuesto computacional definido.

\subsection{Métricas algorítmicas}

\begin{table}[h]
\centering
\begin{tabular}{lll}
\toprule
\textbf{Métrica} & \textbf{Definición} & \textbf{Uso en HCP} \\
\midrule
Costo computacional & Tiempo de ejecución u operaciones realizadas & Comparar algoritmos bajo las mismas condiciones \\
Iteraciones & Número de ciclos completos del algoritmo & Medir convergencia \\
Evaluaciones de fitness & Número de soluciones evaluadas & Métrica justa entre distintos algoritmos \\
Tasa de éxito & Porcentaje de instancias resueltas & Medir efectividad global \\
Tiempo hasta la solución & Tiempo necesario para encontrar el primer ciclo válido & Medir rapidez práctica \\
Estabilidad & Variabilidad entre ejecuciones & Evaluar sensibilidad a la aleatoriedad \\
\bottomrule
\end{tabular}
\end{table}

\subsection{Métricas estadísticas}

Aunque precisión, recall y F1-score provienen de problemas de clasificación, pueden adaptarse para evaluar sistemas que deciden si un grafo es hamiltoniano o no.

La precisión se define como:

\[
\text{Precisión}=\frac{TP}{TP+FP}
\]

El recall o sensibilidad se define como:

\[
\text{Recall}=\frac{TP}{TP+FN}
\]

El F1-score se define como:

\[
F1=\frac{2 \cdot \text{Precisión} \cdot \text{Recall}}{\text{Precisión}+\text{Recall}}
\]

Donde:

\begin{itemize}
    \item \(TP\): verdaderos positivos.
    \item \(FP\): falsos positivos.
    \item \(FN\): falsos negativos.
\end{itemize}

En el contexto del HCP:

\begin{itemize}
    \item Un verdadero positivo es un grafo hamiltoniano correctamente identificado.
    \item Un falso positivo es un grafo no hamiltoniano clasificado incorrectamente como hamiltoniano.
    \item Un falso negativo es un grafo hamiltoniano que el algoritmo no logró resolver.
\end{itemize}

\section{Métricas del FHCP Challenge}

El Flinders Hamiltonian Cycle Project (FHCP) define el HCP como el problema de determinar si un grafo de \(N\) vértices contiene un ciclo simple de longitud \(N\). \cite{fhcp}

\subsection{Características del benchmark}

\begin{itemize}
    \item El FHCP Challenge Set contiene 1{,}001 instancias del HCP. \cite{fhcpset}
    \item Los tamaños de los grafos van desde 66 hasta 9{,}528 vértices. \cite{fhcpset}
    \item El tamaño promedio es ligeramente superior a 3{,}000 vértices. \cite{fhcpset}
    \item Las instancias fueron diseñadas para ser difíciles debido a su estructura, no únicamente por su tamaño. \cite{haythorpe}
    \item Todas las instancias tienen al menos un ciclo hamiltoniano conocido. \cite{haythorpe}
\end{itemize}

\subsection{Formato de los archivos}

Los archivos del FHCP utilizan el formato HCP, estructurado en tres secciones:

\begin{itemize}
    \item Encabezado.
    \item Lista de aristas.
    \item Pie de archivo.
\end{itemize}

El campo \texttt{DIMENSION} indica el número de vértices. El campo \texttt{EDGE\_DATA\_FORMAT : EDGE\_LIST} indica que las conexiones se representan mediante una lista de aristas. \cite{fhcpset}

Los archivos de solución o \textit{tour} contienen la secuencia de vértices en el orden en que son visitados por el ciclo hamiltoniano. \cite{fhcpset}

\subsection{Criterios de comparación}

Para realizar una comparación justa entre algoritmos, se deben fijar:

\begin{itemize}
    \item El mismo subconjunto de instancias.
    \item El mismo límite de tiempo o número de evaluaciones.
    \item Varias ejecuciones independientes por instancia.
    \item La misma definición de éxito.
    \item La verificación de que el ciclo encontrado sea válido.
    \item Métricas de tiempo, evaluaciones, tasa de éxito y desviación estándar.
\end{itemize}

\section{Algoritmos metaheurísticos}

\subsection{Recocido simulado (SA)}

El recocido simulado imita el proceso de enfriamiento controlado de un material. Comienza con una temperatura alta, acepta soluciones peores con cierta probabilidad y reduce gradualmente la temperatura para concentrarse en soluciones de alta calidad.

La probabilidad de aceptación se define como:

\[
P(\text{aceptar})=
\begin{cases}
1, & \Delta f \leq 0 \\
e^{-\Delta f/T}, & \Delta f>0
\end{cases}
\]

donde \(\Delta f\) es el empeoramiento de la función objetivo y \(T\) es la temperatura actual. \cite{kirkpatrick}

\subsection{Ant Colony Optimization (ACO)}

ACO se inspira en el comportamiento colectivo de las hormigas. Cada hormiga construye una solución paso a paso, y las aristas utilizadas en buenas soluciones reciben más feromona.

La probabilidad de seleccionar el siguiente vértice puede expresarse como:

\[
p_{ij}=
\frac{[\tau_{ij}]^\alpha[\eta_{ij}]^\beta}
{\sum_{k\in \mathcal{N}_i}[\tau_{ik}]^\alpha[\eta_{ik}]^\beta}
\]

donde:

\begin{itemize}
    \item \(\tau_{ij}\) representa la feromona.
    \item \(\eta_{ij}\) representa la heurística local.
    \item \(\alpha\) y \(\beta\) controlan la importancia de cada componente.
    \item \(\mathcal{N}_i\) contiene los vecinos válidos no visitados.
\end{itemize}

\subsection{Bee Colony}

El algoritmo de colonia de abejas artificiales divide la colonia en abejas exploradoras, empleadas y observadoras. Las abejas empleadas explotan soluciones actuales, las observadoras seleccionan soluciones prometedoras y las exploradoras buscan nuevas regiones del espacio.

Para el HCP, cada fuente de alimento puede representarse como una permutación de vértices. Los movimientos locales pueden incluir intercambios, inversiones o reparaciones de subciclos. \cite{karaboga}

\subsection{Evolución diferencial (DE)}

La evolución diferencial genera nuevos candidatos combinando diferencias entre individuos de la población. Su operador clásico es:

\[
v_i=x_{r1}+F(x_{r2}-x_{r3})
\]

donde \(x_{r1}\), \(x_{r2}\) y \(x_{r3}\) son individuos distintos y \(F\) es el factor de escala.

Para problemas combinatorios como el HCP, se requiere una adaptación de la codificación, ya que las permutaciones no pueden restarse directamente. Una alternativa es utilizar \textit{random keys}, donde cada individuo se representa mediante valores continuos que posteriormente se ordenan para formar una permutación. \cite{storn}

\subsection{Particle Swarm Optimization (PSO)}

PSO modela un enjambre de partículas que ajustan su posición según su mejor experiencia personal y la mejor experiencia colectiva.

La actualización de velocidad se define como:

\[
v_i^{t+1}=wv_i^t+c_1r_1(p_i-x_i^t)+c_2r_2(g-x_i^t)
\]

La actualización de posición es:

\[
x_i^{t+1}=x_i^t+v_i^{t+1}
\]

donde:

\begin{itemize}
    \item \(w\) es el coeficiente de inercia.
    \item \(c_1\) y \(c_2\) son coeficientes de aprendizaje.
    \item \(r_1\) y \(r_2\) son valores aleatorios.
    \item \(p_i\) es la mejor posición personal.
    \item \(g\) es la mejor posición global.
\end{itemize}

Para el HCP, PSO requiere una codificación discreta, como \textit{random keys}, ya que una permutación no puede representarse directamente mediante valores continuos. \cite{kennedy}

\section{Propuesta de diseño experimental}

Para implementar y comparar los algoritmos, se recomienda seguir los siguientes pasos:

\begin{enumerate}
    \item Definir el HCP como un problema de decisión sobre grafos no dirigidos.
    \item Utilizar permutaciones cíclicas como representación principal.
    \item Validar la factibilidad de cada solución antes de evaluarla.
    \item Definir una función de aptitud que penalice aristas inválidas, subciclos y vértices con grado incorrecto.
    \item Implementar SA, ACO, ABC, DE y PSO bajo la misma interfaz de evaluación.
    \item Utilizar el FHCP Challenge Set como benchmark.
    \item Medir tasa de éxito, tiempo, número de evaluaciones, iteraciones y desviación estándar.
    \item Verificar cada ciclo encontrado antes de reportarlo como solución válida.
    \item Comparar los algoritmos bajo el mismo presupuesto computacional.
\end{enumerate}

La fortaleza del diseño está en separar tres capas:

\begin{itemize}
    \item Representación.
    \item Búsqueda.
    \item Evaluación.
\end{itemize}

Esto permite comparar los cinco algoritmos bajo condiciones equivalentes y obtener resultados reproducibles sobre el FHCP Challenge Set. \cite{haythorpe}

\begin{thebibliography}{9}

\bibitem{west}
West, D. B. (2000). \textit{Introduction to Graph Theory} (2nd ed.). Prentice Hall.

\bibitem{fhcp}
Flinders Hamiltonian Cycle Project. (2020). \textit{Home}. Recuperado de \url{https://sites.flinders.edu.au/flinders-hamiltonian-cycle-project/}

\bibitem{fhcpset}
Flinders Hamiltonian Cycle Project. (2020). \textit{FHCP Challenge Set}. Recuperado de \url{https://sites.flinders.edu.au/flinders-hamiltonian-cycle-project/fhcp-challenge-set/}

\bibitem{haythorpe}
Haythorpe, M. (2019). “FHCP Challenge Set: The First Set of Structurally Difficult Instances of the Hamiltonian Cycle Problem.” \textit{arXiv:1902.10352}.

\bibitem{russell}
Russell, S., \& Norvig, P. (2021). \textit{Artificial Intelligence: A Modern Approach} (4th ed.). Pearson.

\bibitem{kirkpatrick}
Kirkpatrick, S., Gelatt, C. D., \& Vecchi, M. P. (1983). “Optimization by Simulated Annealing.” \textit{Science}, 220(4598), 671–680.

\bibitem{dorigo}
Dorigo, M., Maniezzo, V., \& Colorni, A. (1996). “Ant System: Optimization by a Colony of Cooperating Agents.” \textit{IEEE Transactions on Systems, Man, and Cybernetics, Part B}, 26(1), 29–41.

\bibitem{karaboga}
Karaboga, D., \& Basturk, B. (2007). “A powerful and efficient algorithm for numerical function optimization: artificial bee colony (ABC) algorithm.” \textit{Journal of Global Optimization}, 39(3), 459–471.

\bibitem{storn}
Storn, R., \& Price, K. (1997). “Differential Evolution: A Simple and Efficient Heuristic for Global Optimization over Continuous Spaces.” \textit{Journal of Global Optimization}, 11(4), 341–359.

\bibitem{kennedy}
Kennedy, J., \& Eberhart, R. (1995). “Particle swarm optimization.” \textit{Proceedings of ICNN’95 – International Conference on Neural Networks}, 1942–1948.

\end{thebibliography}

\end{document}


  EOF
